\documentclass{article}
\usepackage[UTF8]{ctex}
\usepackage{amsmath,amssymb,amsfonts}
\usepackage{enumitem}
\usepackage{geometry}
\geometry{a4paper, margin=1in}
% === 错题本排版优化（不参与题目标识） ===
\setlength{\parindent}{0pt}
\setlength{\parskip}{0.35em plus 0.10em minus 0.08em}
\setlength{\emergencystretch}{3em}
\setlist{itemsep=0.25em, topsep=0.40em, parsep=0pt}
\allowdisplaybreaks[2]
% === 错题本排版优化结束 ===
\begin{document}
\section*{华东师范大学2026年数学分析真题}
\section*{38．华东师范大学2026年数学分析试卷}

\subsection{题目1}
1．计算 \(\sum_{n=1}^{\infty} \frac{2 n-1}{2^{n}}=\) \(\_\_\_\_\)．

\subsection{题目2}
2．计算 \(\lim _{x \rightarrow 0^{+}} \sqrt[x]{\sec \sqrt{2 x}}=\)

\subsection{题目3}
3．计算 \(\int_{0}^{1} \int_{0}^{\sqrt{1-x^{2}}} \int_{\sqrt{x^{2}+y^{2}}}^{\sqrt{2-x^{2}-y^{2}}} x y \mathrm{~d} z \mathrm{~d} y \mathrm{~d} x=\)

\subsection{题目4}
4．设 \(|a| \neq 1\) ，计算 \(\int_{0}^{\frac{\pi}{2}} \frac{\arctan (a \tan x)}{\tan x} \mathrm{~d} x=\) \(\_\_\_\_\)．

\subsection{题目5}
5．求和函数 \(\sum_{n=1}^{\infty} \frac{(-1)^{n-1}}{n(n+1)}\left(x^{2}-1\right)^{2 n}=\) \(\_\_\_\_\)．

\subsection{题目6}
6．计算第二型曲面积分 \(\iint_{\sum} x z^{2} \mathrm{~d} y \mathrm{~d} z+\frac{1}{y} \mathrm{~d} z \mathrm{~d} x+x^{2} z \mathrm{~d} x \mathrm{~d} y=\) \(\_\_\_\_\)，其中 \(\sum\) 为抛物面 \(y=x^{2}+z^{2}\)被平面 \(y=1, y=2\) 所截部分，取外侧。

\subsection{题目7}
7．设 \(0<x<2 \pi, a \neq 0\) ，则 \(\frac{1}{2 a}+\sum_{n=1}^{\infty} \frac{a \cos n x-n \sin n x}{n^{2}+a^{2}}=\) \(\_\_\_\_\)．

\subsection{题目8}
8．（15 分）解答如下问题：

\begin{itemize}
  \item[（1）] 证明：对任意正整数 \(n\) ，方程 \(x^{3}+3 x+\frac{1}{n}=0\) 有且仅有一个实数解．
  \item[（2）] 设 \(x_{n}, n=1,2, \cdots\) 为（1）中方程的实根，证明：数列 \(\left\{x_{n}\right\}\) 收敛．
\end{itemize}

\subsection{题目9}
9．（15 分）设 \(\left\{f_{n}(x)\right\}\) 为定义在 \([a, b]\) 上的函数列，对任意 \(x \in[a, b]\) ，数列 \(\left\{f_{n}(x)\right\}\) 有界，且存在 \(K>0\)使得对任意 \(x, y \in[a, b]\) 和 \(n \in \mathbb{N}_{+}\)，都有 \(\left|f_{n}(x)-f_{n}(y)\right| \leq K|x-y|\) 。证明：函数列 \(\left\{f_{n}(x)\right\}\) 必存在一致收敛的子列。

\subsection{题目10}
10．（16 分）判断下列命题是否正确，若正确给出证明，若错误给出反例说明．

\begin{itemize}
  \item[（1）] 若 \(f(x)\) 是 \([a, b]\) 上的凸函数，则 \(f(x)\) 在 \([a, b]\) 上连续．
  \item[（2）] 若 \(f(x)\) 是 \((a, b)\) 上的凸函数，且 \(f(x)\) 在 \((a, b)\) 上可微，则 \(f^{\prime}(x)\) 在 \((a, b)\) 上连续．
\end{itemize}

\subsection{题目11}
11．（16 分）设 \(f(x)\) 在 \(\left[0, \frac{\pi}{2}\right]\) 上为正值连续函数，且

\[
I_{1}=\int_{0}^{\frac{\pi}{2}} e^{f(x)} \sin x \ln f(x) \mathrm{d} x=0, I_{2}=\int_{0}^{\frac{\pi}{2}} e^{f(x)} \cos x \ln f(x) \mathrm{d} x=0
\]

证明：

\begin{itemize}
  \item[（1）] 在 \(\left(0, \frac{\pi}{2}\right)\) 上至少存在 \(x_{1}\) ，使得 \(f\left(x_{1}\right)=1\) ．
  \item[（2）] 在 \(\left(0, \frac{\pi}{2}\right)\) 上至少存在 \(x_{2}, x_{3}\) ，使得 \(f\left(x_{2}\right)=f\left(x_{3}\right)\) ．
\end{itemize}

\subsection{题目12}
12．（16 分）设二元函数 \(f(x, y)\) 的偏导数 \(f_{x}, f_{y}, f_{x y}\) 在点 \(P_{0}\left(x_{0}, y_{0}\right)\) 处都连续．证明：\(f_{y x}\left(x_{0}, y_{0}\right)\) 存在，且 \(f_{y x}\left(x_{0}, y_{0}\right)=f_{x y}\left(x_{0}, y_{0}\right)\) ．

\subsection{题目13}
13．（16 分）设 \(\Omega\) 是 \(\mathbb{R}^{3}\) 中的有界单连通区域，其边界 \(S\) 是一个光滑封闭曲面，\(\vec{\pi}\) 是 \(S\) 上动点 \(Q(\xi, \eta, \zeta)\)处的单位外法向量，\(P_{0}\left(x_{0}, y_{0}, z_{0}\right)\) 为不在 \(S\) 上的一点，\(r=\sqrt{\left(\xi-x_{0}\right)^{2}+\left(\eta-y_{0}\right)^{2}+\left(\zeta-z_{0}\right)^{2}}\) ，证明： \(\iiint_{\Omega} \frac{1}{r} \mathrm{~d} \xi \mathrm{~d} \eta \mathrm{~d} \zeta=\frac{1}{2} \iint_{S} \cos \angle(\overrightarrow{P Q}, \vec{n}) \mathrm{d} S\) ，其中 \(\angle(\overrightarrow{P Q}, \vec{n})\) 表示向量 \(\overrightarrow{P Q}\) 和 \(\vec{n}\) 的夹角。
\end{document}
